\section{Надграждаемост чрез нови конектори}
\label{sec:verification_extensibility}

Една от централните претенции на разработката е, че библиотеката е надграждаема без компромис с нейните гаранции по време на компилация. Това изисква не просто наличие на тестове, а ясно дефиниран договор на конектора (конектор договор), който позволява нов свързващ компонент да бъде добавен без промяна в програмния интерфейс за заявки (заявка API) и без разпадане на типовата дисциплина, ограниченията на възможностите и архитектурните инварианти.

В предходните глави беше показано, че библиотеката вече работи с релационни и нерелационни компоненти, както и че асинхронният слой може да бъде реализиран по различен начин според платформено-специфичните възможности на конкретния драйвер. Настоящата глава формализира как това е възможно и как трябва да се добавя следващ свързващ компонент. С други думи, тук надграждаемостта престава да бъде общо качество на дизайна и се превръща в методология с конкретни изисквания, сигнатури (сигнатури), декларации на възможности, асинхронни задължения (асинхронни задължения) и критерии за завършеност.

\subsection{Формален договор на конекторния слой}

Конекторът се описва чрез тип \code{DB} и специализация \code{connector\_trait<DB>}. Това е структурен, а не наследствен договор. Няма общ виртуален интерфейс \code{IConnector}. Причината е принципна: различните компоненти имат различни възможности и ограничения, които трябва да бъдат видими по време на компилация. Виртуалната йерархия би изравнила твърде много различия и би прехвърлила част от проверките към времето за изпълнение.

Минималният работещ конектор трябва да осигури \code{wire\_type}, \code{cursor\_type} и подходящи входни точки (входни точки) за \code{execute(...)}. Допълнителните способности се означават чрез labelи за възможности. Така договорът остава едновременно строг и разширяем. В тази конструкция няма ``магически'' регистрационен слой: компилаторът вижда специализацията директно и може да провери дали тя удовлетворява концепта (концепт) \code{is\_connector}.

\begin{figure}[H]
\centering
\begin{tikzpicture}[node distance=1.35cm and 1.2cm]
    \node[ormaccent] (dbtype) {\textbf{Нов тип свързващ компонент}\\обвивка на ресурс + жизнен цикъл};
    \node[ormblock, below=of dbtype] (trait) {\code{connector\_trait<DB>}\\задължителни сигнатури};
    \node[ormblock, below left=of trait] (caps) {Декларации\\на възможности};
    \node[ormblock, below right=of trait] (exec) {\code{execute(...)}\\път за материализация};
    \node[ormblock, below=of $(caps)!0.5!(exec)$] (tests) {Модулна + интеграционна\\верификация};

    \draw[ormline] (dbtype) -- (trait);
    \draw[ormline] (trait) -- (caps);
    \draw[ormline] (trait) -- (exec);
    \draw[ormline] (caps) -- (tests);
    \draw[ormline] (exec) -- (tests);
\end{tikzpicture}
\caption{Архитектурен работен процес (workflow) за добавяне на нов конектор}
\label{fig:connector_workflow}
\end{figure}

Фигура~\ref{fig:connector_workflow} показва, че добавянето на компонент не е едноетапно действие. То започва с ясно дефинирана обвивка на ресурс (обвивка на ресурс), продължава с договорена характеристика (договор на характеристика), профил на възможности и материализация на изпълнението, и завършва едва след тестова верификация. Това е съществено за дипломната работа, защото поставя надграждаемостта в дисциплиниран инженерен процес.

\subsection{Задължителни елементи на минимален конектор}

На практика един свързващ компонент се счита за минимално интегриран, ако:

\begin{enumerate}[label=\arabic*.]
    \item съществува тип, който представлява обвивка на връзката на компонента;
    \item съществува специализация \code{connector\_trait<DB>} за този тип;
    \item специализацията дефинира \code{template <typename T> struct wire\_type};
    \item специализацията дефинира \code{cursor\_type};
    \item специализацията може да изпълнява поне базовите ОРС форми на заявки, които твърди, че поддържа.
\end{enumerate}

Това означава, че ``конектор'' не е просто рендерираща функция, а пълно описание на начина, по който C++ типовете и междинното представяне преминават към конкретната целева система. Ако липсва дори един от тези елементи, добавеният компонент не е архитектурно пълноценен, независимо дали част от заявките могат да бъдат изкуствено изпълнени.

\begin{table}[H]
\centering
\small
\caption{Задължителни и опционални елементи на конектора}
\label{tab:connector_requirements}
\begin{tabularx}{\textwidth}{L{4.7cm}Y}
\toprule
\textbf{Елемент} & \textbf{Роля} \\
\midrule
Тип \code{DB} & обвивка над манипулатор на компонента, сесия или имитираща връзка \\[0.5em]
\code{connector\_trait<DB>} & централна специализация на договора \\[0.5em]
\code{template <typename T>\newline struct\ wire\_type} & преобразува C++ тип към мрежово (wire) представяне \\[0.5em]
\code{cursor\_type} & описва начина на обхождане на резултатите или техния жизнен цикъл \\[0.5em]
Претоварвания на \code{execute(...)} & материализират междинното представяне към поведение на компонента \\[0.5em]
Етикети за възможности & ограничават допустимите операции и политиките за жизнен цикъл \\[0.5em]
\code{begin}/\code{commit}/\code{rollback} & необходими при поддръжка на транзакции \\
\code{ddl\_for(...)} & необходимо при интеграция, съобразена със схемата (съобразен със схемата), и \code{migrate<DB>} \\
\code{render\_constexpr(...)} & опционална кука (hook) за материализация по време на компилация за специализирани сценарии \\
\bottomrule
\end{tabularx}
\end{table}

\subsection{Типизиране, сигнатури и именуване}

Договорът на конектора е строг и по отношение на типовете. \code{wire\_type<T>::type} трябва да е стабилна функция по време на компилация на \code{T}. \code{cursor\_type} трябва да моделира реалистично итериране на резултатите или поне ясно да обозначава имитиращо поведение. Претоварванията на \code{execute(...)} трябва да приемат \code{DB\&} като първи аргумент и да връщат \code{orm::result<...>} за \code{SELECT} форми или \code{result<std::tuple<>>} за операции за запис.

Именуването на конекторите в хранилището показва последователна конвенция. ``Локални'' или удобни за модулно тестване (удобни за модулно тестване) конектори използват имена като \code{SQLiteDB}, \code{PostgreSQLDB}, \code{MySQLDB}, \code{MongoDB}, \code{RedisDB}, \code{Neo4jDB} и \code{CassandraDB}. Реалните варианти на живо се разграничават експлицитно чрез суфикса \code{LiveDB}. Това не е дребен стилов избор, а част от яснотата на договора: още от името е ясно дали даден тип е адаптер за имитация (адаптер за имитация), локална/в паметта обвивка или реална целева система, подкрепена от драйвер.

Липсата на наследяване също трябва да се разглежда като формално правило, а не като детайл от реализацията. Нов конектор не трябва да наследява ``обща ОРС база'', защото това би смесило специфичната за компонента логика с полиморфизъм по време на изпълнение. Правилният механизъм за разширение е специализацията на \code{connector\_trait<DB>}.

\subsection{Механизъм на възможностите и ограничения по време на компилация}

Етикетите за възможности са основният начин библиотеката да различава компонентите по допустими операции. Ако \code{connector\_trait<DB>} не декларира \code{supports\_joins}, тогава опит за заявка, базирана на съединяване, е архитектурно неправилен и следва да бъде отхвърлен още при компилация. Същото важи за транзакции, агрегация и конкурентно изпълнение.

От гледна точка на верификацията това е значимо, защото коректността вече не зависи само от тестовете. Част от договора е проверим директно от компилатора. Тестовете тогава не доказват съществуването на договора, а че конкретните специализации го прилагат последователно. Именно по тази логика \path{tests/модулни/test_mongodb_конектор.cpp} и \path{tests/модулни/test_redis_конектор.cpp} валидират отсъствието на \code{supports\_joins}, \code{supports\_transactions} и \code{supports\_aggregation} при съответните конектори.

Тази негативна верификация е също толкова важна, колкото и позитивната. Да докажеш, че компонент \emph{не} поддържа дадена операция, е толкова съществено, колкото да докажеш, че поддържа друга.

\subsection{Минимален и експлоатационен скелет на конектора}

В хранилището \code{MockDB} и \code{MockMigrateDB} са особено ценни, защото играят ролята на архетипове на конектори. Първият показва експлоатационния договор за изпълнение на заявки, а вторият --- как същият договор може да бъде разширен към миграционни сценарии, съобразени със схемата.

\begin{lstlisting}[language=C++,caption={Извадка от \code{connector\_trait<MockDB>}},label={lst:mockdb_connector}]
template <>
struct конектор_trait<MockDB>
{
    using supports_joins              = void;
    using supports_transactions       = void;
    using supports_aggregation        = void;
    using supports_upsert             = void;
    using supports_bulk_insert        = void;
    using supports_concurrent_execute = void;

    template <typename T>
    struct wire_type { using type = T; };

    struct cursor_type
    {
        [[nodiscard]] bool has_next() const noexcept { return false; }
    };

    template <typename Response, typename Joins, typename Wheres,
              typename Limits, typename Groups, typename Orders>
    static auto execute(MockDB& db,
        select_заявка<Response, Joins, Wheres, Limits, Groups, Orders> q)
        -> result<projected_type<Response>, Response>;
};
\end{lstlisting}

Listing~\ref{lst:mockdb_connector} е показателен, защото събира в кратка форма почти всички задължителни елементи: профил на възможностите, \code{wire\_type}, \code{cursor\_type} и специфични за заявките сигнатури за \code{execute(...)}. Освен това конекторът използва отделни претоварвания за \code{select\_query}, \code{insert\_query}, \code{update\_query} и \code{delete\_query}, което прави съпоставянето с поведението на целевата система експлицитно.

\begin{lstlisting}[language=C++,caption={Извадка от разширението \code{MockMigrateDB}, съобразено със схемата},label={lst:mockmigratedb_connector}]
template <>
struct конектор_trait<MockMigrateDB>
{
    using supports_joins              = void;
    using supports_transactions       = void;
    using supports_aggregation        = void;
    using supports_concurrent_execute = void;

    template <typename T>
    struct wire_type { using type = T; };

    template <typename... Args>
    static auto execute(MockMigrateDB& db, Args&&... args)
    {
        return конектор_trait<MockDB>::execute(
            static_cast<MockDB&>(db), std::forward<Args>(args)...);
    }

    [[nodiscard]] static std::string ddl_for(const create_table_op& op);
    [[nodiscard]] static std::string ddl_for(const add_column_op& op);
};
\end{lstlisting}

Listing~\ref{lst:mockmigratedb_connector} показва как експлоатационният скелет на конектора може да бъде надграден. Тук логиката за изпълнение на заявките се делегира към \code{MockDB}, но се добавят DDL куки и куки за транзакции. Това е добър пример за модел на прогресивна зрялост на конектора.

\subsection{Транзакции, нишкова безопасност и подготвени заявки}

Пълноценният договор на целевата система не се изчерпва с единично изпълнение на \code{SELECT}. Ако конекторът декларира \code{supports\_transactions}, от него се очаква да поддържа \code{begin}, \code{commit} и \code{rollback}, така че \code{transaction\_guard<DB>} да бъде валиден. Ако декларира \code{supports\_concurrent\_execute}, той трябва да може да участва безопасно в \code{connection\_pool<DB, N>}.

Слоят за подготвени заявки също поставя изискване за стабилност на жизнения цикъл на \code{DB\&} и за консистентно поведение при свързване при многократно изпълнение. Следователно зрелият конектор трябва да бъде оценяван не само по това дали рендира правилен синтаксис, а и по експлоатационния договор около ресурси, транзакции и повторна употреба. Това е особено важно за компоненти, които имат по-сложен жизнен цикъл на драйвера от \code{MockDB}.

\subsection{Миграции и интеграция, съобразена със схемата}

Конектор, който участва в \code{migrate<DB>}, трябва да предложи и точки за разширение на DDL. В хранилището това е демонстрирано чрез \code{MockMigrateDB}. Този пример е особено полезен, защото показва как върху вече съществуваща целева система за изпълнение може да се надгради поведение, съобразено със схемата. По този начин библиотеката формализира различни нива на завършеност на конектора.

Тук е важно да се отбележи, че поддръжката, съобразена със схемата, не е задължителна за всяка целева система. Някои конектори могат да бъдат напълно адекватни за анализ на изпълнението на заявките, без да поддържат автоматизирани миграции. Но ако конекторът претендира за по-пълна ОРС интеграция, наличието на \code{ddl\_for(...)} куки и миграционни тестове вече става задължително.

\begin{figure}[H]
\centering
\begin{tikzpicture}[node distance=1.35cm and 1.15cm]
    \node[ormaccent] (lvl1) {Ниво 1\\минимален конектор};
    \node[ormblock, below=of lvl1] (lvl2) {Ниво 2\\експлоатационен конектор\\изпълнение на заявки + договор за резултати};
    \node[ormblock, below=of lvl2] (lvl3) {Ниво 3\\съобразен с транзакции/конкурентност};
    \node[ormblock, below=of lvl3] (lvl4) {Ниво 4\\съобразен със схемата конектор\\DDL + куки за миграция};
    \node[ormnote, right=1.5cm of lvl3, minimum width=4.7cm] (note)
        {\code{MockDB} покрива експлоатационна зрялост, а
         \code{MockMigrateDB} показва преход към зрялост, съобразена със схемата.};

    \draw[ormline] (lvl1) -- (lvl2);
    \draw[ormline] (lvl2) -- (lvl3);
    \draw[ormline] (lvl3) -- (lvl4);
\end{tikzpicture}
\caption{Стълба на зрелостта за нов конектор}
\label{fig:connector_maturity_ladder}
\end{figure}

Фигура~\ref{fig:connector_maturity_ladder} предлага практичен модел за оценка. Тя е полезна и за самото бъдещо развитие на библиотеката, защото позволява новите компоненти да бъдат позиционирани не бинарно като ``има/няма конектор'', а според нивото на тяхната интеграция.

\subsection{Методика за добавяне на нов свързващ компонент}

Добавянето на нов компонент следва ясна последователност. Първо се дефинира типът на целевата система и неговата обвивка на ресурс. След това се реализира \code{connector\_trait<DB>} със задължителните елементи. Следват декларациите на възможности и правилата за рендиране или изпълнение на междинното представяне. Накрая се добавят модулни тестове за договора и, при възможност, интеграционни тестове на живо за реалния драйвер.

Архитектурно правилният ред е важен: не се започва от заобиколни решения по време на изпълнение, а от формализиране на границите по време на компилация и на логиката на съпоставяне. Ако този ред бъде нарушен, резултатът обикновено е конектор, който ``работи'' за единичен идеален сценарий (идеален сценарий), но не е съвместим с общата архитектура.

\begin{table}[H]
\centering
\small
\caption{Препоръчителна матрица за верификация за нов конектор}
\label{tab:connector_verification_matrix}
\begin{tabularx}{\textwidth}{L{3.0cm}L{3.3cm}Y}
\toprule
\textbf{Пласт на верификация} & \textbf{Какво трябва да се докаже} & \textbf{Примерен корелат в хранилището} \\
\midrule
Съответствие с концепт & \code{is\_connector<DB>} е удовлетворен & \path{tests/модулни/test_mongodb_конектор.cpp} \\
Честност на възможностите & декларираните и недекларираните labelи за възможности са правилни & \path{tests/модулни/test_redis_конектор.cpp}, \path{tests/интеграционни/test_sqlite.cpp} \\
Коректност на изпълнението & междинното представяне се материализира коректно & \path{tests/интеграционни/test_mockdb.cpp} и тестовете на живо \\
Коректност на ресурси/жизнен цикъл & куките за транзакции и конкурентност се държат коректно & \path{tests/модулни/test_нишка_safety.cpp} \\
Коректност, съобразена със схемата & DDL куките и миграционните разлики работят коректно & \path{tests/модулни/test_schema_migration.cpp} \\
\bottomrule
\end{tabularx}
\end{table}

\subsection{Критерии за „напълно имплементиран и функциониращ“ конектор}

На базата на кода и тестовете в хранилището могат да се формулират следните критерии:

\begin{enumerate}[label=\arabic*.]
    \item конекторът удовлетворява договора \code{is\_connector};
    \item декларациите на възможности съответстват на реалното поведение;
    \item конвенциите за именуване и типизиране са съвместими с останалите конектори;
    \item ресурсният жизнен цикъл е коректен и проверим;
    \item поддържаните ОРС операции са покрити от модулни тестове;
    \item при наличен драйвер на живо има интеграционни тестове срещу реална база данни;
    \item при заявена поддръжка, съобразена със схемата, са налични DDL куки и миграционни тестове.
\end{enumerate}

Тези критерии са по-строги от обикновено ``има работещ адаптер'', защото изискват консистентност между архитектурно намерение, договорен програмен интерфейс и верифицирано поведение. Именно това различава непланираната (непланирана) интеграция от конектор, който наистина принадлежи към библиотеката.

\subsection{Верификация чрез тестове}

Тестовият слой в хранилището следва структурата на архитектурата. Модулните тестове за отделните конектори валидират съответствието с \code{is\_connector}, декларациите на възможности и основното рендиране. \path{tests/модулни/test_нишка_safety.cpp} проверява \code{connection\_pool}, \code{thread\_local\_db} и \code{transaction\_guard}. \path{tests/модулни/test_schema_migration.cpp} валидира логиката за разлики в миграциите. \path{tests/интеграционни/test_mockdb.cpp} и интеграционните тестове на живо за отделни системи доказват, че абстракцията работи и при реално изпълнение.

Следователно верификацията не е концентрирана само на едно място. Тя е разпределена така, че всяко архитектурно твърдение да има поне един пряк тестов корелат. Това превръща разширяемостта на конекторите от декларативна цел в проверима инженерна практика.

\subsection{Оценка на наличните конектори}

\code{SQLiteDB} може да се разглежда като референтна зряла реализация за релационна целева система. \code{PostgreSQLDB}, \code{MySQLDB}, \code{MongoDB}, \code{RedisDB}, \code{Neo4jDB} и \code{CassandraDB} вече демонстрират важни части от договора и имат различна степен на покритие на живо. \code{MockDB} и \code{MockMigrateDB} играят особена роля като верификационни конектори: те не са само тестови двойници, а средство за формализиране на правилата на архитектурата.

Оттук следва и по-общият извод на настоящата глава: надграждаемостта на библиотеката не е добавена впоследствие, а е заложена в самия начин, по който е дефиниран конекторният слой. Именно това позволява разработката да бъде продължавана с нови свързващи компоненти, без да се нарушава формалният характер на ОРС архитектурата по време на компилация.
